\begin{tomtat}
	\begin{minipage}[t]{0.3\textwidth}
		\vspace{0pt}
		\begin{tomtat}
			\begin{center}
				\textbf{ĐỀ 08}
			\end{center}
			
		\end{tomtat}
	\end{minipage}
	\hfill
	\begin{minipage}[t]{0.5\textwidth}
		\vspace{0pt}
		\begin{tomtat}
			\begin{center}
				\includegraphics[scale=0.08]{images/8}
			\end{center}
		\end{tomtat}
	\end{minipage}
\end{tomtat}
\textbf{I. PHẦN TRẮC NGHIỆM}
\begin{ex}
	Giải phương trình sau: $(2x+6)(x-3)=0$.
	\choice{$-3$ và $3$}{$-6$ và $3$}{$3$}{$-3$}
	\loigiai{
		Ta có phương trình tích $(2x+6)(x-3)=0$.
		Điều này xảy ra khi một trong hai thừa số bằng 0:
		\begin{itemize}
			\item $2x+6=0 \Rightarrow 2x=-6 \Rightarrow x=-3$
			\item $x-3=0 \Rightarrow x=3$
		\end{itemize}
		Vậy nghiệm của phương trình là $x=-3$ và $x=3$.
		\\ \textbf{ĐÁP ÁN A}
	}
\end{ex}

\begin{ex}
	Giải bất phương trình sau: $5x-12>6x-15$.
	\choice{$x>3$}{$x>-3$}{$x<3$}{$x<-3$}
	\loigiai{
		Ta giải bất phương trình:
		\begin{align*}
			5x-12 &> 6x-15 \\
			-12+15 &> 6x-5x \\
			3 &> x \\
			x &< 3
		\end{align*}
		Vậy nghiệm của bất phương trình là $x<3$.
		\\ \textbf{ĐÁP ÁN C}
	}
\end{ex}

\begin{ex}
	Giá trị biểu thức $\frac{1}{2-\sqrt{3}}+\frac{1}{2+\sqrt{3}}$
	\choice{$x=-4$}{$x=4$}{$x=2\sqrt{3}$}{$x=-2\sqrt{3}$} 
	% Lưu ý: Các lựa chọn có vẻ không phù hợp với câu hỏi tính giá trị biểu thức. 
	% Giả sử đây là lỗi và ta sẽ tính giá trị.
	\loigiai{
		Ta quy đồng mẫu số:
		\begin{align*}
			A &= \frac{1(2+\sqrt{3})}{(2-\sqrt{3})(2+\sqrt{3})} + \frac{1(2-\sqrt{3})}{(2+\sqrt{3})(2-\sqrt{3})} \\
			&= \frac{2+\sqrt{3}}{2^2 - (\sqrt{3})^2} + \frac{2-\sqrt{3}}{2^2 - (\sqrt{3})^2} \\
			&= \frac{2+\sqrt{3}}{4 - 3} + \frac{2-\sqrt{3}}{4 - 3} \\
			&= \frac{2+\sqrt{3}}{1} + \frac{2-\sqrt{3}}{1} \\
			&= 2+\sqrt{3} + 2-\sqrt{3} \\
			&= 4
		\end{align*}
		Giá trị của biểu thức là 4. Nếu phải chọn một đáp án từ các lựa chọn đã cho, thì đáp án B ($x=4$) có giá trị số bằng 4.
		\\ \textbf{ĐÁP ÁN B} (Giả sử câu hỏi muốn tìm giá trị số là 4)
	}
\end{ex}

\begin{ex}
	Tính thể tích V của hình cầu có bán kính $R=3\text{cm}$ là
	\choice{$35\pi \text{ cm}^3$}{$36\pi \text{ cm}^3$}{$36,5\pi \text{ cm}^3$}{$37\pi \text{ cm}^3$}
	\loigiai{
		Thể tích của hình cầu được tính bằng công thức: $V = \frac{4}{3}\pi R^3$.
		Với $R=3 \text{ cm}$:
		\begin{align*}
			V &= \frac{4}{3}\pi (3)^3 \\
			&= \frac{4}{3}\pi \cdot 27 \\
			&= 4\pi \cdot \frac{27}{3} \\
			&= 4\pi \cdot 9 \\
			&= 36\pi (\text{cm}^3)
		\end{align*}
		\\ \textbf{ĐÁP ÁN B}
	}
\end{ex}

\begin{ex}
	Một cột đèn có bóng trên mặt đất dài 8,5m. Các tia nắng mặt trời tạo với mặt đất một góc xấp xỉ bằng $52^\circ$. Tính chiều cao của cột đèn (làm tròn đến hàng phần mười).
	\choice{$10,9\text{cm}$}{$10,8\text{m}$}{$10,8\text{cm}$}{$10,9\text{m}$}
	\loigiai{
		Gọi $h$ là chiều cao của cột đèn (cạnh đối) và $b = 8,5 \text{ m}$ là độ dài bóng trên mặt đất (cạnh kề).
		Góc tạo bởi tia nắng mặt trời và mặt đất là $\alpha = 52^\circ$.
		Ta có mối quan hệ: $\tan \alpha = \frac{\text{cạnh đối}}{\text{cạnh kề}} = \frac{h}{b}$.
		Do đó, $h = b \cdot \tan \alpha$.
		$$ h = 8,5 \cdot \tan(52^\circ) $$
		Sử dụng máy tính: $\tan(52^\circ) \approx 1,2799$.
		$$ h \approx 8,5 \cdot 1,2799 \approx 10,879 \text{ m} $$
		Làm tròn đến hàng phần mười, ta được $h \approx 10,9 \text{ m}$.
		\\ \textbf{ĐÁP ÁN D}
	}
\end{ex}

\begin{ex}
	Tìm điều kiện xác định của: $\frac{2}{\sqrt{x^2-4x+4}}$.
	\choice{$x>2$}{$x<2$}{$x \ne 2$}{$x=2$}
	\loigiai{
		Biểu thức xác định khi mẫu số khác 0 và biểu thức dưới dấu căn bậc hai lớn hơn hoặc bằng 0.
		Tuy nhiên, vì căn bậc hai nằm ở mẫu, nên biểu thức dưới dấu căn phải lớn hơn 0 (không được bằng 0).
		Ta có: $\sqrt{x^2-4x+4} = \sqrt{(x-2)^2} = |x-2|$.
		Để biểu thức xác định, ta cần $|x-2| > 0$.
		Điều này xảy ra khi $x-2 \ne 0$, tức là $x \ne 2$.
		\\ \textbf{ĐÁP ÁN C}
	}
\end{ex}

\begin{ex}
	Ba bạn Việt, Nam, Lào được xếp ngẫu nhiên ngồi trên một hàng ghế có ba chỗ ngồi. Xác suất của các biến cố: Việt không ngồi ngoài cùng bên phải là:
	\choice{$\frac{1}{3}$}{$\frac{2}{3}$}{$1$}{$\frac{4}{3}$}
	\loigiai{
		Gọi 3 vị trí ngồi là 1 (trái), 2 (giữa), 3 (phải).
		Ký hiệu Việt (V), Nam (N), Lào (L).
		Tập hợp tất cả các kết quả có thể xảy ra khi xếp 3 bạn vào 3 chỗ ngồi (không gian mẫu $\Omega$):
		\begin{itemize}
			\item (V, N, L), (V, L, N)
			\item (N, V, L), (N, L, V)
			\item (L, V, N), (L, N, V)
		\end{itemize}
		Số phần tử của không gian mẫu là $n(\Omega) = 3! = 6$.
		
		Tập hợp các kết quả thuận lợi cho biến cố A: "Việt không ngồi ngoài cùng bên phải".
		Điều này có nghĩa là Việt ngồi ở vị trí 1 (trái) hoặc vị trí 2 (giữa).
		\begin{itemize}
			\item Việt ngồi vị trí 1: (V, N, L), (V, L, N) - 2 trường hợp.
			\item Việt ngồi vị trí 2: (N, V, L), (L, V, N) - 2 trường hợp.
		\end{itemize}
		Vậy, các kết quả thuận lợi cho biến cố là: (V, N, L), (V, L, N), (N, V, L), (L, V, N).
		Số phần tử của biến cố thuận lợi là $n(A) = 4$.
		
		Xác suất của biến cố là:
		$$ P(A) = \frac{n(A)}{n(\Omega)} = \frac{4}{6} = \frac{2}{3} $$
		\\ \textbf{ĐÁP ÁN B}
	}
\end{ex}

\begin{ex}
	Gọi $x_1, x_2$ là hai nghiệm của phương trình $x^2-5x+3-2x=0$. Tính $x_1+x_2=?$
	\choice{7}{-5}{5}{-7}
	\loigiai{
		Đầu tiên, ta rút gọn phương trình đã cho:
		$x^2 - 5x + 3 - 2x = 0$
		$x^2 - (5x+2x) + 3 = 0$
		$x^2 - 7x + 3 = 0$.
		Đây là phương trình bậc hai có dạng $ax^2+bx+c=0$ với $a=1, b=-7, c=3$.
		Theo hệ thức Vi-ét, tổng hai nghiệm $x_1+x_2$ được tính bằng công thức:
		$$ x_1+x_2 = \frac{-b}{a} = \frac{-(-7)}{1} = \frac{7}{1} = 7 $$
		\\ \textbf{ĐÁP ÁN A}
	}
\end{ex}
\newpage
\textbf{II. PHẦN TỰ LUẬN}
\begin{ex} % Câu 9
	Giải hệ phương trình sau: $\begin{cases} 4x + 3y = 1 \\ 2x - 3y = 11 \end{cases}$.
	\loigiai{
		Ta có hệ phương trình:
		$$ \begin{cases} 4x + 3y = 1 & (1) \\ 2x - 3y = 11 & (2) \end{cases} $$
		Cộng vế theo vế của phương trình (1) và phương trình (2), ta được:
		\begin{align*}
			(4x + 3y) + (2x - 3y) &= 1 + 11 \\
			6x &= 12 \\
			x &= \frac{12}{6} \\
			x &= 2
		\end{align*}
		Thay $x=2$ vào phương trình (1):
		\begin{align*}
			4(2) + 3y &= 1 \\
			8 + 3y &= 1 \\
			3y &= 1 - 8 \\
			3y &= -7 \\
			y &= -\frac{7}{3}
		\end{align*}
		Vậy nghiệm của hệ phương trình là $(x; y) = \left(2; -\frac{7}{3}\right)$.
	}
\end{ex}

\begin{ex} % Câu 11 (Đánh số lại từ ảnh)
	Rút gọn biểu thức sau: $A = \left(\frac{2}{\sqrt{x}-2} - \frac{\sqrt{x}+4}{x-4}\right) : \frac{\sqrt{x}}{\sqrt{x}+2}$, $x > 0, x \ne 4$.
	\loigiai{
		Với $x > 0, x \ne 4$, ta có:
		\begin{align*}
			A &= \left(\frac{2}{\sqrt{x}-2} - \frac{\sqrt{x}+4}{(\sqrt{x}-2)(\sqrt{x}+2)}\right) : \frac{\sqrt{x}}{\sqrt{x}+2} \\
			&= \left(\frac{2(\sqrt{x}+2) - (\sqrt{x}+4)}{(\sqrt{x}-2)(\sqrt{x}+2)}\right) \cdot \frac{\sqrt{x}+2}{\sqrt{x}} \\
			&= \frac{2\sqrt{x}+4 - \sqrt{x}-4}{(\sqrt{x}-2)(\sqrt{x}+2)} \cdot \frac{\sqrt{x}+2}{\sqrt{x}} \\
			&= \frac{\sqrt{x}}{(\sqrt{x}-2)(\sqrt{x}+2)} \cdot \frac{\sqrt{x}+2}{\sqrt{x}}
		\end{align*}
		Vì $x>0$ nên $\sqrt{x} \ne 0$. Vì $x \ne 4$ nên $\sqrt{x}+2 \ne 0$ và $\sqrt{x}-2 \ne 0$. Ta có thể rút gọn $\sqrt{x}$ và $\sqrt{x}+2$.
		\begin{align*}
			A &= \frac{1}{\sqrt{x}-2}
		\end{align*}
		Vậy $A = \frac{1}{\sqrt{x}-2}$.
	}
\end{ex}

\begin{ex}
	Cho phương trình $x^2 - 2x - 4 = 0$, không giải phương trình, hãy tính giá trị biểu thức sau $T = x_1^2 - x_2(2-x_2) - 2x_1$.
	\loigiai{
		Phương trình đã cho là $x^2 - 2x - 4 = 0$.
		Đây là phương trình bậc hai có dạng $ax^2+bx+c=0$ với $a=1, b=-2, c=-4$.
		Ta tính biệt thức $\Delta = b^2-4ac = (-2)^2 - 4(1)(-4) = 4 + 16 = 20$.
		Vì $\Delta = 20 > 0$ nên phương trình có hai nghiệm phân biệt $x_1, x_2$.
		
		Theo hệ thức Vi-ét, ta có:
		$$ \begin{cases} 
			x_1+x_2 = \dfrac{-b}{a} = \dfrac{-(-2)}{1} = 2 \\ 
			x_1x_2 = \dfrac{c}{a} = \dfrac{-4}{1} = -4 
		\end{cases} $$
		Biến đổi biểu thức $T$:
		\begin{align*}
			T &= x_1^2 - x_2(2-x_2) - 2x_1 \\
			&= x_1^2 - 2x_2 + x_2^2 - 2x_1 \\
			&= (x_1^2 - 2x_1) + (x_2^2 - 2x_2)
		\end{align*}
		Vì $x_1$ là nghiệm của phương trình $x^2 - 2x - 4 = 0$, nên $x_1^2 - 2x_1 - 4 = 0 \Rightarrow x_1^2 - 2x_1 = 4$.
		Tương tự, vì $x_2$ là nghiệm của phương trình $x^2 - 2x - 4 = 0$, nên $x_2^2 - 2x_2 - 4 = 0 \Rightarrow x_2^2 - 2x_2 = 4$.
		Thay vào biểu thức $T$:
		\begin{align*}
			T &= 4 + 4 \\
			&= 8
		\end{align*}
		Vậy giá trị của biểu thức $T$ là 8.
	}
\end{ex}

\begin{ex} % Câu 12. a) (Đánh số lại từ ảnh)
	Bạn An chơi trò chơi gieo đồng thời hai con xúc xắc cân đối đồng chất I và II để xem qua các lần gieo xúc xắc mình có thể gieo nhiều nhất được bao nhiêu điểm. Tính xác suất của các biến cố A: "Có đúng 1 con xúc xắc xuất hiện mặt 6 chấm".
	\loigiai{
		Tập hợp tất cả các kết quả có thể xảy ra khi gieo đồng thời hai con xúc xắc (không gian mẫu $\Omega$) là các cặp (Xúc xắc I, Xúc xắc II): \\
		(1,1), (1,2), (1,3), (1,4), (1,5), (1,6),
		(2,1), (2,2), (2,3), (2,4), (2,5), (2,6),
		(3,1), (3,2), (3,3), (3,4), (3,5), (3,6),
		(4,1), (4,2), (4,3), (4,4), (4,5), (4,6),
		(5,1), (5,2), (5,3), (5,4), (5,5), (5,6),
		(6,1), (6,2), (6,3), (6,4), (6,5), (6,6).\\
		Số phần tử của không gian mẫu là $36$.
		
		Tập hợp các kết quả thuận lợi cho biến cố A: "Có đúng 1 con xúc xắc xuất hiện mặt 6 chấm":\\
		(1,6), (2,6), (3,6), (4,6), (5,6),
		(6,1), (6,2), (6,3), (6,4), (6,5).\\
		Số kết quả thuận lợi cho biến cố A là $n(A) = 10$.
		
		Xác suất của biến cố A là:
		$$ P(A) = \frac{n(A)}{n(\Omega)} = \frac{10}{36} = \frac{5}{18} $$
	}
\end{ex}

\begin{ex}
	Một hội trường có 300 ghế ngồi được xếp đều thành từng dãy (số ghế trên mỗi dãy bằng nhau). Nếu kê thêm 2 dãy và mỗi dãy xếp thêm 1 ghế thì trong hội trường có tất cả 351 ghế. Hỏi lúc đầu hội trường có bao nhiêu dãy ghế.
	\loigiai{
		Gọi số dãy ghế lúc đầu là $x$ (dãy).
		Điều kiện: $x \in \mathbb{N}^*$.
		
		Số ghế trên mỗi dãy lúc đầu là: $\frac{300}{x}$ (ghế/dãy).
		
		Sau khi thay đổi:
		Số dãy ghế mới là $x+2$ (dãy).
		Số ghế trên mỗi dãy mới là $\frac{300}{x} + 1$ (ghế/dãy).
		Tổng số ghế mới trong hội trường là $(x+2)\left(\frac{300}{x} + 1\right)$.
		
		Theo đề bài, tổng số ghế mới là 351, nên ta có phương trình:
		$$ (x+2)\left(\frac{300}{x} + 1\right) = 351 $$
		$$ (x+2)\left(\frac{300+x}{x}\right) = 351 $$
		$$ (x+2)(300+x) = 351x $$
		$$ 300x + x^2 + 600 + 2x = 351x $$
		$$ x^2 + 302x + 600 = 351x $$
		$$ x^2 + 302x - 351x + 600 = 0 $$
		$$ x^2 - 49x + 600 = 0 $$
		Tính biệt thức $\Delta = b^2 - 4ac = (-49)^2 - 4(1)(600) = 2401 - 2400 = 1$.
		$\sqrt{\Delta} = \sqrt{1} = 1$.
		Phương trình có hai nghiệm:
		$$ x_1 = \frac{-(-49) + 1}{2(1)} = \frac{49 + 1}{2} = \frac{50}{2} = 25 $$
		$$ x_2 = \frac{-(-49) - 1}{2(1)} = \frac{49 - 1}{2} = \frac{48}{2} = 24 $$
		Cả hai giá trị $x=25$ và $x=24$ đều là số nguyên dương.
		Kiểm tra lại:
		\begin{itemize}
			\item Nếu $x=25$ dãy:
			Số ghế mỗi dãy ban đầu: $300/25 = 12$ ghế.
			Số dãy mới: $25+2 = 27$ dãy.
			Số ghế mỗi dãy mới: $12+1 = 13$ ghế.
			Tổng số ghế mới: $27 \times 13 = 351$ ghế (đúng).
			\item Nếu $x=24$ dãy:
			Số ghế mỗi dãy ban đầu: $300/24 = 12,5$ ghế (loại vì số ghế trên mỗi dãy phải là số nguyên).
		\end{itemize}
		Vậy lúc đầu hội trường có 25 dãy ghế.
	}
\end{ex}

\begin{ex} 
	Cho tam giác $ABC$ nhọn. Đường tròn $(O)$ đường kính $BC$ cắt $AB, AC$ lần lượt tại $E$ và $D$; $BD$ cắt $CE$ tại $H$, $AH$ cắt $BC$ tại $I$. Từ $A$ kẻ tiếp tuyến $AM, AN$ của đường tròn $(O)$ ($M, N$ là tiếp điểm).
	\begin{enumerate}
		\item[a)] Chứng minh tứ giác $AEHD$ nội tiếp.
		\item[b)] Chứng minh $AB \cdot BE = BI \cdot BC$ từ đó suy ra $AB \cdot BE + AC \cdot CD = BC^2$.
	\end{enumerate}
	\loigiai{
		\begin{center}
			\begin{tikzpicture}[line join = round, line cap=round,>=stealth,font=\footnotesize,scale=1]
				% Def
				\def\r{3} % r=3cm
				
				% Coordinates
				\coordinate (O) at (0,0);
				\coordinate (B) at (-\r,0);
				\coordinate (C) at (\r,0);
				\coordinate (A) at (-1,5);
				
				\coordinate (E) at ($(A)!(C)!(B)$); % E la hinh chieu cua C len AB
				\coordinate (D) at ($(A)!(B)!(C)$); % D la hinh chieu cua B len AC
				% I la hinh chieu cua A len BC. Vi BC nam tren truc Ox, I se co hoanh do cua A va tung do 0.
				\coordinate (I) at ($(A|-O)!(A)!(B)$); % Cach khac: \coordinate (I) at (-1,0);
				% Hoac: \coordinate (I) at ($(B)!(A)!(C)$);
				
				% H la giao diem cua BD va CE (gia su H la truc tam)
				\path[name path=BD_line] (B) -- (D);
				\path[name path=CE_line] (C) -- (E);
				\coordinate[name intersections={of=BD_line and CE_line, by=H}];
				
				% Dinh nghia duong tron duong kinh AO
				\coordinate (K1) at ($(A)!0.5!(O)$); % K1 la trung diem AO, tam duong tron duong kinh AO
				\path let \p1 = ($(A)-(O)$), \n1 = {veclen(\x1,\y1)/2} in 
				[name path=circle_diam_AO, draw=none] (K1) circle (\n1);
				
				% Dinh nghia duong tron duong kinh BC (chinh la duong tron (O,r))
				\path [name path=circle_diam_BC, draw=none] (O) circle (\r);
				
				% Tim giao diem M, N cua hai duong tron tren
				% Su dung mot ten tam thoi cho N de tranh xung dot neu N da duoc dinh nghia
				\coordinate[name intersections={of=circle_diam_AO and circle_diam_BC, by={M,N_intersect}}];
				% Doi ten N_intersect thanh N de su dung trong cac lenh ve sau
				\coordinate (N) at (N_intersect);
				
				% Draw
				\draw (O) circle (\r); % Ve duong tron (O,r)
				
				\draw (A) -- (B) -- (C) -- cycle; % Ve A-B-C-A
				\draw (A) -- (I); % Ve A-I
				\draw (B) -- (D); % Ve B-D
				\draw (C) -- (E); % Ve C-E
				
				\draw[dashed] (O) -- (M); % Ve net dut O-M
				\draw[dashed] (O) -- (N); % Ve net dut O-N
				
				\draw (A) -- (M) -- (I) -- (N) -- cycle; % Ve A-M-I-N-A
				
				% Ky hieu goc vuong
				% Can \usetikzlibrary{angles, quotes} o preamble
				\pic[draw, angle radius=2.5mm]{right angle = A--M--O}; % Goc AMO vuong
				\pic[draw, angle radius=2.5mm]{right angle = A--N--O}; % Goc ANO vuong
				
				% Goc AEH: E la chan duong cao tu C. Neu H la truc tam, H nam tren CE.
				% Ky hieu goc vuong tai E do CE vuong goc AB
				\pic[draw, angle radius=2.5mm]{right angle = C--E--A}; % Goc CEA vuong (hoac CEB)
				% Goc ADH: D la chan duong cao tu B. Neu H la truc tam, H nam tren BD.
				% Ky hieu goc vuong tai D do BD vuong goc AC
				\pic[draw, angle radius=2.5mm]{right angle = B--D--A}; % Goc BDA vuong (hoac BDC)
				
				% Fill 
				\fill (O) circle (1.5pt);
				\fill (A) circle (1.5pt);
				\fill (B) circle (1.5pt);
				\fill (C) circle (1.5pt);
				\fill (E) circle (1.5pt);
				\fill (D) circle (1.5pt);
				\fill (I) circle (1.5pt);
				\fill (H) circle (1.5pt);
				\fill (M) circle (1.5pt);
				\fill (N) circle (1.5pt);
				% \fill (K1) circle (1pt); % Trung diem K1 khong can hien thi
				
				% Nodes	
				\node[below left] at (O) {$O$};
				\node[above] at (A) {$A$};
				\node[below left] at (B) {$B$};
				\node[below right] at (C) {$C$};
				\node[above left] at (E) {$E$};
				\node[below] at (D) {$D$};
				\node[below] at (I) {$I$};
				\node[left] at (H) {$H$}; % Vi tri H co the can dieu chinh
				\node[left] at (M) {$M$}; % Vi tri M co the can dieu chinh
				\node[right] at (N) {$N$};  % Vi tri N co the can dieu chinh
			\end{tikzpicture}
			
		\end{center}
		}
\end{ex}

\begin{ex} % Câu 14
	\begin{enumerate}
		\item[a)] Một xí nghiệp sản xuất được 180 sản phẩm loại I và 180 sản phẩm loại II trong thời gian 5 giờ. Mỗi giờ sản xuất được số sản phẩm loại I nhiều hơn số sản phẩm loại II là 30 sản phẩm. Hỏi mỗi giờ xí nghiệp sản xuất được bao nhiêu sản phẩm mỗi loại?
		\item[b)] Tìm giá trị nhỏ nhất của biểu thức: $A = 4x + \frac{1}{4x} - \frac{4\sqrt{x}+3}{x+1} + 2027$ với $x > 0$.
	\end{enumerate}
\end{ex}


\setcounterref{ex}{0}